% Chapter content to be written.
In this chapter, we \todo{something}

\todo{at the end, give an overview of primary results here}
\todo{include the stuff from original info which gives some background on the various explicit construction problems such as rigidity. Say something brief like: many of these problems are famou and have recieved intense scrutiny, we give a very brief overview of some of the work here.}
\section{Explicit Constructions and Circuit Lower Bounds are Total Search Problems}
\todo{this is brief: give the key perspective here, that we can think of explicit constructions as total search problems and proving circuit lower bounds for exponential time classes as algorithms for total search problems}
\todo{perhaps, remove the general definition of ``explicit construction search problems'' from the preliminaries and port it over to this section}
\todo{ok but unsure how much of this should be in the intro vs this chapter}
\todo{give an example of this perspective: Impagliazzo/nisan/wigderson is a polynomial time reduction from one explicit construction problem to another}

\section{The Range Avoidance Search Problem}
\begin{enumerate}
    \item Remind reader (again) the definition of the problem
    \item Give basic upper/lower bounds for the problem
    \item relate it to classical hardness/randomness theory
\end{enumerate}
\section{Reductions to Range Avoidance}
We start with \todo{sparse string encodings, remote point}
\subsection{Hard Functions, Random Strings, and Pseudorandom Generators}
The first and most important reduction we cover is the construction of truth tables of boolean functions which require exponential circuit complexity. In this setting, the input length $N$ of our explicit construction problem is a power of 2 and represents the truth table of a boolean function:
\todo{modify later terminology to match this}
\begin{definition}
Let $s: \mathbb{N} \to \mathbb{N}$ be some ``size bound.'' Define the search problem {\sc $s(n)$- Hard} as follows: given $1^{2^n}$, output the truth table of a boolean function $f: \B^n \to \B$ such that $x$ is not computed by any circuit of size at most $s(n)$.
\end{definition}
Shannon famously showed that, for $s(n) \leq \frac{2^n}{n}$, this problem is \emph{total} for $n$ sufficiently large. We show below that for $s(n) \leq \frac{2^n}{2n}$, this problem reduces in polynomial time to Avoid. Recall that the input to this problem is $1^{2^n}$, which has length $2^n$, so when we say our reduction runs in ``polynomial time'' we mean $\poly(2^n)$ time.

\begin{theorem}\label{thm:truth-tables-in-apepp}
{\sc $\frac{2^n}{2n}$-Hard} reduces in polynomial time to Avoid.
\end{theorem}
\begin{proof}
This proof directly follows Shannon's aforementioned classical argument. We construct an instance of Avoid in the form of a circuit $C$ which maps an encoding of a circuit to its corresponding truth table. $C$ interprets its input as a circuit on $n$ bits (using an encoding of circuits to be described below), evaluates it on every possible input to generate a $2^{n}$ bit truth table, and outputs the truth table.

We now describe the encoding of circuits used by $C$, which will guarantee that any length-$2^n$ truth table $f$ which has a circuit of size $\frac{2^n}{2n}$ is in the range of $C$. Given a circuit of size $s \geq n $ on $n$ inputs computing $f$, for each of its $s$ gates we can use 2 bits to encode whether it is an $\land, \lor$, or $\lnot$ gate, and an additional $2\log (s+n) \leq 2\log s + O(1)$ bits to specify its inputs. We can then use an additional $\log s$ bits to specify which gate is the terminal output gate. Overall this requires $2s\log s + O(s)$ bits. It is clear that from such an encoding, $C$ can efficiently decode the represented circuit and evaluate it on all possible input values. For $s \leq \frac{2^n}{2n}$, we have that $2 s \log s + O(1) < 2^n$ for $n$ sufficiently large. So $C$ is indeed a valid instance of Avoid, and any string outside the range of $C$ is a solution to {\sc $\frac{2^n}{2n}$-Hard}. It is also clear from the above description that $C$ can be constructed in $poly(2^n)$ time.

\end{proof}

Following the original publication of this work, it was shown in \cite{CHR} that the above argument could be sharpened using a more involved circuit encoding procedure:
\begin{theorem}\cite{CHR}
    {\sc $\frac{2^n}{n}$-Hard} reduces in polynomial time to Avoid.
\end{theorem}

A related problem was studied in \cite{total}, referred to there as ``{\sc Complexity}.'' Rather then an explicit construction problem with sparse input, in this problem you are given as input one truth table $f:\B^n \to \B$, and asked to produce another truth table $g$ such that $g$ requires large circuits, even with access to $f$-oracle gates. More formally:

\begin{definition}\label{def:complexity}
The problem {\sc Complexity} is defined as follows: given an $f: \B^n \to \B$, find another $g: \B^n \to \B $ such that $g$ requires $f$-oracle circuits of size $\frac{2^n}{n^2}$.
\end{definition}
In \cite{total} the following is shown:
\begin{theorem}
{\sc Complexity} reduces in polynomial time to {\sc Avoid}.
\end{theorem}
\todo{mention Ilango's proof that this is hard, if the oracle is a bit longer}

Finally, we discuss another important generalization of Theorem~\ref{thm:truth-tables-in-apepp} due to subsequent work \cite{RSW}. An important contribution of \cite{RSW} was to emphasize the importance of ``parameterizing'' Avoid instances in two key ways. For an Avoid instance $C: \B^n \to \B^m$, we may classify it according to (1) the \emph{stretch} of the instance, which denotes the quantitative relationship between $m$ and $n$, and (2) the \emph{circuit complexity} $C$, e.g. whether it is a constant depth boolean formula, an unbounded depth formula, or a fully general boolean circuit. 
When reducing the construction of hard truth tables to Avoid, there is a natural tradeoff between the quantitative strength of the lower bound and stretch of Avoid instance: if we wish to produce $f: \B^n \to \B$ hard for circuits of size $s$, we can reduce to an Avoid instance which stretches $\poly(s(n)) $ bits to $ 2^n$ bits.

More interestingly, \cite{RSW} show that if we wish to produce explicit truth tables hard for some class of circuits $\cl{C}$ (e.g. formulas, branching programs, circuits), we can reduce this problem to an Avoid instance $C$ where each output bit of $C$ is in the \emph{same class $\cl{C}$}. This will hold (using the same proof as in Theorem~\ref{thm:truth-tables-in-apepp}) provided that the ``truth table generator'' mapping circuits in the class $\cl{C}$ to their truth tables is computable in the same class $\cl{C}$. This is sometimes referred to as a circuit class having the ``universal property,'' and many important classes of restricted circuits are known to have it (see \cite{RSW} for a further discussion). As a consequence we have:

\begin{theorem}[\cite{RSW}]\label{thm:strong-tt-to-avoid}
\sloppy Let $\cl{C}$ be a circuit class with the universal property, for example $\cl{C} \in \{\mathsf{AC^0},\mathsf{ACC^0}, \mathsf{TC^0}, \mathsf{NC^1}, \mathsf{L}\}$. Then producing $f: \B^n \to \B$ requiring $\cl{C}$-circuits of size $s(n)$ is polynomial time reducible to solve a $\cl{C}$-Avoid instance stretching $\poly(s(n))$ bits to $ {2^n}$ bits.

As a corollary, if Avoid for $\cl{C}$-circuits is solvable in polynomial time (resp. with an $\mathsf{NP}$ oracle) then $\mathsf{E}$ (resp. $\mathsf{E^{NP}}$) requires $\cl{C}$-circuits of size $2^{\Omega(n)}$.
\end{theorem}
\todo{about the corollary: when do we first bring up this connection between circuit lower bounds and explicit construction of hard truth tables? need to have that here, and mention it right below the first Shannon-hard reduction}

% Note that, for circuit classes such as $\mathsf{TC^0}$ and $ \mathsf{NC^1}$, it an important open problem to show that $\mathsf{E^{NP}}$ does not have $\poly(n)$ size circuits from the class. By the above connection, this means that we would already achieve a breakthrough from, for example, a $2^{O(n)}$ time $\mathsf{NP}$ oracle algorithm for $\mathsf{TC^0}\dash\avoid$ in the stretch regime $n^{\omega(1)} \mapsto 2^n$, i.e. when the number of outputs is almost exponential in the number of inputs.

Next, we turn our attention to the explicit construction of pseudorandom generators.
\begin{definition}
We will say that a sequence $R = (x_1, \ldots, x_m)$ of $n$-bit strings is a pseudorandom generator if, for all $n$-input circuits of size $n$:
\[
|Pr_{i \sim m}[C(x_i)=1] - Pr_{x \sim \B^n}[C(x)=1]| \leq 1/n
\]
\end{definition}

Standard applications of the probabilistic method show that such pseudorandom generators exist of size polynomial in $n$. In particular, for some polynomial $p(n)$ we may define the following explicit construction problem:
\begin{definition}
{\sc $p(n)$-Prg} is the following search problem: given $1^n$, output a pseudorandom generator $R = (x_1, \ldots, x_m)$, $x_i \in \B^n$, $m = p(n)$.
\end{definition}

A polynomial time algorithm for {\sc $p(n)$ Prg} would suffice to derandomize \cl{BPP} \cite{NW} \todo{link this to earlier discussion}. We will show here that the explicit construction of a PRG of size $n^6$ is reducible to Avoid.
As noted in \todo{somewhere}, the results of Impagliazzo and Wigderson \cite{IW} imply that {\sc Prg} reduces directly to {\sc $2^{\epsilon n}$-Hard} for any constant $\epsilon > 0$, so a reduction of PRG construction to Avoid follows from Theorem~\ref{thm:truth-tables-in-apepp}. However, we provide here a much simpler direct proof that PRG construction reduces to Avoid, relying only on Yao's next bit predictor lemma, and neither the nearly disjoint subsets construction of Nisan and Wigderson \cite{NW} nor the rather involved worst-case to average-case reductions of Impagliazzo and Wigderson \cite{IW}. Together with our completeness result in Section~\ref{sec:completeness} \todo{fix section reference}, this gives an alternative, self-contained proof that worst-case-hard truth tables can be used to construct pseudorandom generators (although it yields a weaker result, as our derandomization will require an \cl{NP} oracle).

The basic idea of the proof is as follows: if $R$ fails to be a pseudorandom generator for circuits of size $n$, then applying Yao's lemma there exists a next bit predictor circuit $P$ of size $O(n)$ which can predict the $i^{th}$ bit of most strings in $R$ from the previous $i-1$ bits. This will allow us to compress $R$ by omitting the $i^{th}$ bits of all the strings in $R$ and replacing them with the predictor $P$ and some advice bits which will correct the errors made by $P$.

\begin{theorem}\label{thm:prg}
{\sc $n^6$-Prg} reduces in polynomial time to {\sc Empty}
\end{theorem}

\begin{proof}
We start by describing the Avoid instance $C$. $C$ interprets its input as representing a list $R^- = (z^1, \ldots, z^{n^6})$ of $n^6$ strings of length $n-1$, a circuit $P$ of size $cn$ (for a fixed universal constant $c$ to be defined later) with $n-1$ input bits and one output bit, a $\log n$-bit index $j \in [n]$, and a string $S$ encoding an $n^6$-bit string with weight at most $\frac{n^6}{2} - n^4$. Given these, $C$ feeds each $n-1$-bit string $z^i \in R^-$ through $P$, then inserts the output as an extra bit at the $j^{th}$ position of $z^i$ to obtain an $n$-bit string $x^i$. In the final step, $C$ decodes the $n^6$-bit string represented by $S$, and flips the $j^{th}$ bit of $x^i$ if and only if the $i^{th}$ position of $S$ has a 1. $C$ then outputs $R = (x^1, \ldots, x^{n^6})$.

It is clear that $C$ can be implemented as a circuit of size polynomial in $n$, and further that $C$ can be constructed in polynomial time given $1^n$. We now show there are fewer inputs than outputs. Note that the input size is equal to the number of bits needed to specify $R^-, P, j, S$, which is $n^7 - n^6 + \Tilde{O}(n) + \log n + bits(S)$, where $bits(S)$ is the number of bits needed to specify $S$. Since $S$ an $n^6$-bit string with weight at most $\frac{n^6}{2} - n^4 = n^6(\frac{1}{2} - \frac{1}{n^2})$, we can apply Lemma~\ref{lem:sparse-enc}\todo{make sure we are using correct reference} with $\epsilon=\frac{1}{n^2}$ to give an encoding for $S$ using $n^6(1-\frac{1}{n^4}) + \log (n^6) = n^6 - n^2 + 6\log n$ bits. Thus the overall number of bits of input is at most $n^7 - n^6 + \Tilde{O}(n) + \log n + n^6 - n^2 + 6\log n = n^7 + \Tilde{O}(n) + 7\log n - n^2$ which is strictly less then the $n^7$ bits of output (for sufficiently large $n$). So this is indeed a polynomial time reduction to a valid instance of Avoid. It remains to show that any string outside the range of $C$ is a pseudorandom generator.

Let $R = (x^1, \ldots, x^{n^6})$ be a sequence of $n$-bit strings which fails to be a pseudorandom generator; we aim to show that $R$ lies in the range of $C$. By definition, there exists a circuit $D$ of size $n$ such that:
\[
|Pr_{i \sim [n^6]}[D(x^i)=1] - Pr_{x\sim \B^n}[D(x)=1]| > 1/n
\]

For $x \in \B^n$ and $j \in [n]$, let $x_{-j}$ be the $n-1$-bit string obtained by deleting the $j^{th}$ bit of $x$. By Yao's next bit predictor lemma \cite{Yao} \cite{Vadhan}, the previous inequality implies the existence of an index $j \in [n]$ and a circuit $P: \B^{n-1} \rightarrow \B$ of size $cn$ for some fixed universal constant $c$, such that
\[
Pr_{i \sim [n^6]}[P(x_{-j}) = x_j] > \frac{1}{2} + \frac{1}{n^2}
\]
Since $R$ has size $n^6$, this implies that $P$ correctly guesses the $j^{th}$ bit of $x^i$ from the other $n-1$ bits for at least $n^6(\frac{1}{2} + \frac{1}{n^2}) = \frac{n^6}{2} + n^4$ of indices $i \leq n^6$. So if we let $S$ be the $n^6$-bit string with a 1 at the indices $i$ where $P$ guesses the wrong value of $x^i_j$, we see that the weight of $S$ is at most $\frac{n^6}{2} - n^4$ as required. Taking $R^- = (x^1_{-j}, \ldots, x^{n^6}_{-j})$, we have that from $R^-,P,j,S$, our Avoid instance will successfully reconstruct $R$. Thus, any string outside the range of $C$ is a pseudorandom generator.
\end{proof}

Finally, we turn to the construction of strings of high time-bounded Kolmogorov complexity:
\begin{definition}
Let $U$ be any fixed Turing machine, and let $T: \mathbb{N} \rightarrow \mathbb{N}$ be a time bound. For a string $x$, $\cl{K}_U^T(x)$ denotes the length of the shortest string $y$ such that $U$ outputs $x$ on input $y$ in $\leq T(|x|)$ steps. 

For a Turing machine $U$, time bound $T$, and ``length'' parameter $\ell: \mathbb{N} \to \mathbb{N}$ we define the following explicit construction problem {\sc $(\cl{K}^T_U, \ell)$-Random}: given $1^n$, output an $n$ bit string $x$ such that $K_U^T(x) \geq \ell(n)$. In the special case $\ell(n) = n-1$ we omit the parameter $\ell$.
\end{definition}

\begin{theorem}\label{thm:kt-random}
For any fixed Turing machine $U$ and fixed polynomial time bound $T$, {\sc $K_U^T$-Random} reduces in polynomial time to {\sc Empty}.
\end{theorem}
\begin{proof}
We construct a circuit $C$ with $n-1$ input bits and $n$ output bits as follows. If the input is of the form $0^*1y$ ($0^*$ denotes an arbitrary length string of zeros), $C$ simulates $U$ on $y$ for $t(n)$ steps and outputs the result, padded/truncated to length $n$ if it is not already. Otherwise, $C$ outputs $0^n$. It is clear that $C$ can be produced in polynomial time given $1^n$ for any fixed $U,t$. Further, for any $n$-bit string $x$ with $K_U^t(x) \leq n - 2$, there is some $y$ of length $k \leq n-2$ such that $U$ outputs $x$ on input $y$ in $t(n)$ steps, and so $C$ will output $x$ on input $0^{n-k-2}1y$. So any string $x$ outside the range of $C$ must have $K_U^t(x) \geq n - 1$.
\end{proof} 
\subsection{Strongly Explicit Extractors}
Here we describe a technique which allows us to reduce to Avoid the construction of circuits computing functions $f: \B^n \to \B^n$ with certain strong pseudorandom properties. The reductions here run in time $\poly(n)$, and produce a circuit for $f$ of size $\poly(n)$; this is in contrast to the reduction in \todo{hard truth table} which produce $n$-bit boolean functions in time $2^{O(n)}$ given as truth tables.

The reductions here hinge on a simple high-level construction involving $k$-wise independent generators. To explain the method in sufficient generality we require the following definition:
\begin{definition}
Let $\mathcal{A} \subseteq \binom{\B^n}{k}$, $\mathcal{B} \subseteq (\B^{n})^k$ be given. We say that $f: \B^n \to \B^n$ is a $(\mathcal{A},\mathcal{B})$ extractor if whenever $(x_1,\ldots,x_k) \in \mathcal{A}$, $(f(x_1),\ldots, f(x_k)) \notin \mathcal{B}$.
\end{definition}
The idea behind the definition is as follows. We think of $\mathcal{A} \subseteq \{S \subseteq \B^n, |S| = k\}$ as defining a class of \emph{structured} subsets of $k$ inputs, with $|\mathcal{A}| << \binom{2^n}{k}$. For any fixed set of $k$ inputs $(x_1,\ldots,x_k)$ and a function $f: \B^n \to \B^n$, we may think of $(f(x_1),\ldots,f(x_k)) \in \mathcal{B}$ as some unlikely event, i.e. $|\mathcal{B}| << 2^{nk}$. We then have that $f$ is an $(A,B)$ extractor iff $f(S)$ avoids the event $f(S) \in \mathcal{B}$ whenever $S \in \mathcal{U}$. If $\Pr[(y_1,\ldots,y_k) \in \mathcal{B}]$ is very small for a random sequence $y_1,\ldots,y_k \sim \B^n$, then for any fixed $S = \{x_1,\ldots,x_k\} \in \mathcal{U}$ of size $k$, if we choose $f$ uniformly at random we are very likely to have $f(S) = \{f(x_1),\ldots,f(x_k)\} \notin \mathcal{B}$. If $|\mathcal{A}|$ is also sufficiently small than we can union bound over $S \in \mathcal{A}$ and argue that a random $f$ will be an $(\mathcal{A},\mathcal{B})$ extractor with high probability. 

This existence argument can be improved in the following sense: say that instead of choosing the function $f \sim \B^n \to \B^n$ uniformly at random, we sampled it from a $k$-wise independent distribution $\mathcal{F}$: for every fixed $\{x_1,\ldots,x_k\} \in \B^n$ distinct, $(f(x_1),\ldots,f(x_k))$ is distributed uniformly on $(\B^{n})^k$ when $f \sim \mathcal{F}$. Then in this case we have $f(S) \notin \mathcal{B}$ with the same high probability over $f \sim \mathcal{F}$, and the rest of the argument goes through as before. Thus we can argue for the existence of an $(\mathcal{A},\mathcal{B})$ extractor inside of any $k$-wise independent function family $\mathcal{F}$.

The following reduction formalizes this using the standard construction of $k$-wise independent families by univariate polynomials over a finite field. The reduction works provided that the ``smallness'' of $\mathcal{A}, \mathcal{B}$ can also be exhibitted by reductions to Range Avoidance, i.e. we can cover these sets by the range of some explicit length-expanding functions $A: \B^{\ell} \to \binom{\B^n}{k}$, $B: \B^{r} \to \B^{nk}$. The technique here is inspired by a construction of Razborov \cite{razborov-ramsey} showing that strongly explicit Ramsey graphs exist in $\cl{AC^0[2]}$.
\begin{theorem}\label{thm:general-extractor-to-avoid}
Let $A: \B^{\ell} \to \binom{\B^n}{k}$, $B: \B^{r} \to \B^{nk}$ be given as circuits. Provided $\ell + r < nk$, the following problem is reducible in polynomial time to range avoidance: output the description of a polynomial size circuit $f: \B^n \to \B^n$ which is an $(\rng(A), \rng(B))$ extractor. Moreover:
\begin{enumerate}
	\item The circuit $f$ produced by the reduction lies in $\mathsf{AC}^0[2]$.
	\item If $A, B \in \mathsf{NC}^2$ the Range Avoidance instance produced by the reduction also lies in $\mathsf{NC}^2$.
\end{enumerate} 
\end{theorem}
\begin{proof}
We construct an instance of range avoidance $C: \B^{\ell + r} \to \B^{nk}$. Elements of $\B^{nk}$ are interpreted as vectors $\bar{\alpha}=(\alpha_1,\ldots,\alpha_k) \in \mathbb{F}_{2^n}$, and to each $\bar{\alpha}$ we associate the function $f_{\bar{\alpha}}: \B^n\to \B^n$ given by:
\[
f_{\bar{\alpha}}(x) = \sum_{i = 1}^k \alpha_i x^{i-1}
\]
where $\B^n$ and $\mathbb{F}_{2^n}$ are associated in some standard way. We will construct $C$ so that $f_{\alpha}$ fails to be a $(\rng(A),\rng(B))$ extractor, then $\bar{\alpha} \in \rng(C)$ which will give the theorem. The fact that $f_{\bar{\alpha}}$ lies in $\mathsf{AC}^0[2]$ for any choice of $\bar{\alpha}$ follows from results of \cite{finite-field} on the complexity of arithmetic over $\mathbb{F}_{2^n}$.

$C$ is constructed as follows. We interpret the input as encoding $(a,b)$ with $a \in \B^{\ell}$ and $b \in \B^{r}$. Let $(x_1,\ldots,x_k) = A(a)$ and let $(v_1,\ldots, v_k) = B(b)$. We then compute the unique degree $k-1$ polynomial $p \in \mathbb{F}_{2^n}[x]$ such that $p(x_i)=v_i$ for all $i$; this can be accomplished in polynomial time using Gaussian elimination on the corresponding Vandermonde matrix (by assumption $A(a)$ outputs a list of $k$ distinct strings). Finally $C$ outputs the coefficients $(\alpha_1,\ldots,\alpha_k)$ of $p$. Clearly we have $\bar{\alpha} \in \rng(C)$ whenever $f_{\bar{\alpha}}$ fails to be an $(\rng(A), \rng(B))$ extractor. The complexity of $C$ is upper bounded by the complexity of computing $A,B$, together with that of Gaussian elimination over $\mathbb{F}_{2^n}$ which lies in $\mathsf{NC}^2$.
\end{proof}
The above Lemma can be used to reduce the construction of various ``structured seed extractors'' with near-optimal parameters to Range Avoidance. An important case of a structured seed extractor is a 2-source extractor: 
\begin{definition}
We say that a function $f:\{0,1\}^n \times \{0,1\}^n \rightarrow \{0,1\}$ is a $(k,\epsilon)$ 2-source extractor if the following holds: for any two sets $X,Y \subseteq \{0,1\}^n$ of size $2^k$, $|Pr_{x \sim X, y \sim Y}[f(xy) = 1] - \frac{1}{2}| \leq \epsilon$.
\end{definition}
To see that this is a special case of our general $(\mathcal{A}, \mathcal{B})$-extractor paradigm, consider our two-source extractor $f$ to be a function $f: \B^{2n} \to \B^{2n}$, for which we only care about it's first output bit. The family of simple sets $\mathcal{A}$ is the set of ``combinatorial rectangles:'' sets of the form $L \times R$ for some $L, R \subseteq \B^n$, $|L|=|R|=2^k$. Such a combinatorial rectangle has size $2^{2k}$ but can be encoded using $2^{k+1}n$ bits, which yields the appropriate circuit $A$ for the reduction. The ``bad event'' $\mathcal{B}$ is that the first bit of the function $f$ is biased towards 1 or 0 over the rectangle; here we can use the sparse string encodings from \todo{ref} to form the circuit $B$. More formally we have:
\begin{theorem}\label{thm:two-source-to-avoid}
There is a $\poly(n, \epsilon^{-1})$-time reduction from strongly explicit constructing $( \log n + 2 \log (1/\epsilon) + 2, \epsilon)$ 2-source extractors to Avoid. The circuits output by the reduction lie in $\cl{AC^0_2}$.
\end{theorem}
\begin{proof}
Let $m= 2^k$. We require circuits $A,B$ of the following form. $A$ takes $2mn$ bits of input and outputs subsets of $\B^{2n}$ of size $m^2$; the range of $A$ contains all sets of the form $\{ (\ell,r) \mid \ell \in L, r \in R\}$ for some $|L| = |R| = m$, $L, R \subseteq \B^n$. $B$ takes as input strings of length $m^2(2n-1) + m^2 - \epsilon^2 m^2 + O(\log m)$ and outputs elements of $(\B^{2n})^{m^2}$; it's range contains all sequences $x^1, \ldots, x^{m^2} \in \B^{2n}$ such that the bits $x_1^1,\ldots,x_1^{m^2}$ havd bias $>\epsilon$ (the fraction which equal 1 differs from $\frac{1}{2}$ by more than $\epsilon$). $A$ may be constructed in the obvious way (it takes the description of $L,R$ and outputs all cross-concatenations). $B$ may be constructed using \todo{ref}. We may now apply Theorem~\ref{thm:general-extractor-to-avoid} to get a reduction to Avoid, provided $2mn + m^2(2n-1) + m^2 - \epsilon^2 m^2 + O(\log m) < 2nm^2$; for $n$ sufficiently large this is satisfied provided $m \geq \frac{4n}{\epsilon^2}$, in other words $k \geq \log n + 2 \log(1/\epsilon) + 2$.
\end{proof}

Recent advances have yielded 2-source extractors with parameters very close to optimal, but still not meeting the bounds attainable here \cite{li-extractor,extractors-survey}. Other kinds of structured-seed extractors studied in the literature include extractors for varieties and for affine spaces, see \cite{varieties,affine}; similar reductions exist for constructions of each of these using Theorem~\ref{thm:general-extractor-to-avoid}.
\subsection{Miscellanious Reductions}
We discuss here a few more important explicit construction problems which can be shown reducible to Range Avoidance. The first is the problem of \emph{matrix rigidity}, famously introduced by Valiant \cite{Valiant}:

\begin{definition}\cite{Valiant}
We say that $n \times n$ matrix $M$ over $\mathbb{F}_q$ is $(r, s)$ rigid if for any matrix $S \in \mathbb{F}_q^{n \times n}$ with at most $s$ non-zero entries, $M+S$ has rank greater than $r$.
\end{definition}

\begin{definition}
For any $q: \mathbb{N} \rightarrow \mathbb{N}$ such that $q(n)$ is a prime power for all $n$, and let $r, s: \mathbb{N} \to \mathbb{N}$ be given. Then {\sc $(r(n), s(n), q)$-Rigid} is the following search problem: given $1^n$, output an $n \times n$ matrix $M$ over $\mathbb{F}_{q(n)}$ which is $(r(n), s(n))$ rigid.
\end{definition}
% epsilon, q rigid was the problem of rank \epsilon n, sparsity \epsilon n^2
\begin{theorem}\label{thm:rigid-in-apepp}
For any constant $\epsilon>0$ so that $H(\epsilon) + 3\epsilon < 1$, and any efficiently computable $q(n)$ satisfying the above, {\sc $(\epsilon n, \epsilon n^2, q(n))$-Rigid} reduces in polynomial time to Avoid.
\end{theorem}
In particular $\epsilon = \frac{1}{8}$ suffices.
\begin{proof}
Let $M$ be any $n \times n$ matrix over $\mathbb{F}_q$ which is not $(r,s)$ rigid. So there exists an $n \times r$ matrix $L$, an $r \times n$ matrix $R$, and an $n \times n$ matrix $S$ with at most $s$ non-zero entries, such that $M = LR + S$. It is clear that from the descriptions of $L,R,S$ we can efficiently compute $M$.

$L$ and $R$ can each be described explicitly using $nr \log q$ bits. For $S$, we encode it by specifying an $n^2$-bit string $T$ of weight $s$ denoting the entries of $S$ which are nonzero, together with an $s \log q$-bit string $C$ giving the correct values of the nonzero entries. Using this scheme we may assume $s$ is exactly equal to $\epsilon n^2$ since if it is smaller we can add some arbitrary indices and put their unchanged values in $C$. Applying the encoding scheme in Lemma~\ref{lem:sparse-enc} for $T$, overall the number of bits in this encoding is at most $\log\binom{n^2}{s} + (2nr + s)\log q$. Setting $r = \epsilon n$ and $s= \epsilon n^2$ this is at most:
\begin{gather*}
\log\binom{n^2}{\epsilon n^2} + (3\epsilon n^2)\log q
 \leq \\
H(\epsilon)n^2 + (3 \epsilon n^2)\log q \leq \\
(H(\epsilon) + 3\epsilon)n^2\log q
\end{gather*}
By our choice of $\epsilon$ this is bounded above by $(1- \Omega(1))n^2\log q$ bits, whereas an arbitrary matrix in  $\mathbb{F}_q^{n \times n}$ requires $n^2\log q$ bits.
\end{proof}

In section \todo{direc redctuons} of this chapter, we will show that 2 problems, namely the construction of communication matrices outside $\cl{PSPACE^{cc}}$ and data structure problems with near maximal bit-probe complexity, can be reduced directly to the construction of truth tables of circuit complexity $\frac{2^n}{n}$ (or indeed, complexity $s(n)$ for various smaller exponential bounds $s(n)$). Using Theorem~\ref{thm:truth-tables-in-apepp} this yields a reduction of each of these explicit construction problems to Avoid; we will leave the formal definition of these problems until Section \todo{section}. Finally, we include one more important problem which was subsequently shown to reduce to Avoid in \cite{guruswami2022range}:
\begin{definition}[GV Codes]
A matrix \todo{finish this}
\end{definition}
\section{Completeness}
In this section, we exhibit several problems which are \emph{equivalent} to Avoid under suitable reductions. \todo{say more}. 

We start with the following simple completeness result, which uses a more restricted form of reduction than the subsequent reductions: we only need to make \emph{nonadaptive} calls to an $\cl{NP}$ oracle for this reduction, whereas the subsequent reductions will require adaptive calls to an $\cl{NP}$ oracle.
\begin{theorem}
Let $\epsilon > 0$ be any fixed constant, $T(n) = n^{1+\epsilon}$, $\ell(n) = n - n^{\epsilon}$. Let $U$ be any fixed universal Turing machine. Then {\sc $(\cl{K}^T_U, \ell)$-Random} is equivalent to Avoid under many-one $\cl{FP^{NP}_\Vert}$ reductions.
\end{theorem}
\begin{proof}
We have already shown a polynomial time many-one reduction from $(\cl{K}^T_U, \ell)$ to Avoid in the previous section; we now show the reverse direction. Let $C: \B^n \to \B^{n+1}$ be an avoid instance and let $s$ be the circuit size of $C$. For some parameter $k$ to be chosen later, consider the circuit $C^{\otimes k} : \B^{kn} \to \B^{kn + k}$ which is obtained by concatenating $k$ copies of $C$. Let $m = kn + k$. Any string in the range of $C^{\otimes k}$ has $T:= O(m + ks)$-time bounded Kolmogorov complexity at most $s + m + O(\log m)$: we simply supply the description of $C$, together with the length $kn$ input and a $O(\log m)$ size program description of the following program: given a length $m$ string and a description of $C$, break the $m$-bit string into $k$ contiguous chunks of length $n$, applies $C$ to each chunk, and outputs the result. Setting $k\geq (s+ n)^{\frac{2}{\epsilon}}$ we have that $n, s \leq k^{\epsilon/2} \leq m^{\epsilon/2}$, therefore $T \leq m^{1 + \epsilon}$ and $kn \leq m - m^\epsilon$. Now, if $y \in \B^m$ is a string with $T$-time bounded Kolmogorov complexity $> kn$, then, after dividing $y$ into $k$ blocks of $n+1$ bits, we obtain a list of $k$ strings, one of which must be outside the range of $C$. Using $k$ parallel calls to an $\cl{NP}$ oracle we may determine which of these are indeed Avoid solutions for $C$ and output the lexicographically first.
\end{proof}
\todo{do we remove most of the ``emtpy pigeonhole'' terminology in this chapter?}
For our next results, we first require the following reduction which allows us to increase the \emph{stretch} of Avoid instances using an $\cl{NP}$ oracle. This result \todo{is from prior work}
\begin{lemma}\label{lem:swapping-sizes}
There is a ${\cl FP^{NP}}$ many-one reduction from Avoid to itself, which maps instances $C: \B^n \to \B^{n+1}$ to instances $C': \B^n \to \B^{2n}$.
\end{lemma}
\begin{proof}
Let $C: \B^n \to \B^{n+1}$ be given. For a positive integer $i$, we define $C^i: \{0,1\}^n \rightarrow \{0,1\}^{n + i}$ inductively as follows. For the base case, $C^1$ is simply defined as $C$. For $i > 1$, we first compute $C^{i-1}$ on the input $x$ to obtain a string $x'$ of length $n + i - 1$, and then compute $C$ on the first $n$ bits of $x'$ and concatenate the output with the remaining bits of $x'$. Since $C$ replaces the first $n$ bits with $n+1$ bits, if $C^{i-1}$ has output length $n + i-1$ then $C^i$ has output length $n + i$. We claim that for any positive $i$, if we are given an $n+i$-bit string outside the range of $C^i$, we can find an $n$-bit string outside the range of $C$ in $i \cdot poly(|C|)$ time using an $\cl{NP}$ oracle. In particular, we prove by induction on $i$ that this can be accomplished with $i$ calls to an $\cl{NP}$ search solver. For $i=1$ this is trivially the case, since $C^1=C$. Now say $i > 1$. Let $y = xz$ be outside the range of $C^i$ for some $x \in \{0,1\}^{n+1}$, $z \in \{0,1\}^{i-1}$. We use the $\cl{NP}$ search solver to find a preimage of $x$ or report that none exist. If $x$ has no preimage we have solved Avoid for $C$ and are finished. Otherwise, we are given some preimage $x' \in \{0,1\}^n$ for $x$ under $C$. So then $x'z$ must be outside the range of $C^{i-1}$, since if it were not then $xz=y$ would have a preimage under $C^i$. By induction we may use this string $x'z$, and $i-1$ addition $\cl{NP}$ search oracle calls, to find an Avoid solution for $C$. Setting $i = n$ we get the stated result.
\end{proof}

\todo{fix references below}
\begin{theorem}\label{thm:epsilon-hard-reduction}
Let $\epsilon > 0$ be any constant. Then Avoid is $\cl{FP^{NP}}$ many-one reducible to {\sc $2^{\epsilon n}$-Hard}.
\end{theorem}

\begin{proof}
By \todo{ref} we know that Avoid is ${\cl FP^{NP}}$ many-one reducible to the case of instances stretchuing $n$ bits to $2n$ bits. Now, let $C: \B^n \to \B^{2n}$ be such an instance and let $k = 2 \lceil \log |C| \rceil \lceil \frac{1}{\epsilon}\rceil$. Consider the following map $C^*: \{0,1\}^n \rightarrow \{0,1\}^{2^kn}$, defined informally as follows: given a string $x \in \{0,1\}^n$, apply $C$ once to get two $n$-bit strings, then apply $C$ to both of those $n$-bit strings to get four, and continue $k$ times until we have $2^k$ many $n$-bit strings, or equivalently a $2^kn$-bit string. This process is illustrated in Figure~\ref{fig:circuit_tree} below:

\begin{figure}[H]
    \centering
    \includegraphics[width=0.7\textwidth]{circuit_tree.pdf}
    \caption{Extending a map $C: \{0,1\}^n \rightarrow \{0,1\}^{2n}$ to a map $C^*: \{0,1\}^n \rightarrow \{0,1\}^{2^kn}$. Dotted boxes indicate the number of bits along a wire.}
    \label{fig:circuit_tree}
\end{figure}

To define this function more formally, first we define the following maps $L,R: \{0,1\}^{2n} \rightarrow \{0,1\}^n$, where $L$ takes a $2n$-bit string and outputs the first $n$ bits, and $R$ takes a $2n$-bit string and outputs the last $n$ bits. Given a nonempty sequence $\sigma_1, \ldots \sigma_t \in \{L,R\}^*$, let $C^{\sigma}: \{0,1\}^n \rightarrow \{0,1\}^n$ be the function $\sigma_t \circ C \circ \ldots \circ \sigma_1 \circ C$. Now, given a binary string, we can associate it with such a sequence by associating 0 with $L$ and 1 with $R$, and so we will abuse notation and write $C^x$ for a binary string $x$ as shorthand for $C^{\sigma}$ where $\sigma$ is the sequence of $L,R$ associated with the binary string $x$. We are now ready to formally define our function $C^*$. As $C^*$ is a map $\{0,1\}^n \rightarrow \{0,1\}^{2^kn}$, we can think of the output as being clumped into $2^k$ blocks, each containing $n$ bits. Given this terminology, $C^*$ behaves as follows: on input $x$, the $i^{th}$ block of the output of $C^*$ will be $C^{\|i\|}(x)$, where $\|i\|$ denotes the standard representation of $i$ as a $k$-bit binary string.

From here, the proof proceeds in two steps. First, we show that, by setting $m = n 2^k = poly(|C|)$, any solution to {\sc $2^{\epsilon n}$-Hard} on input $1^m$ will be a string that is not in the range of $C^*$. Second, we will show that given a string outside the range of $C^*$, we can find a string outside the range of $C$ using only a polynomial number of calls to an $\cl{NP}$-oracle.

To carry out the first of these steps, we will show that any string in the range of $C^*$, when interpreted as a truth table of length $m = n 2^k$ on $\lceil \log n \rceil + k$ variables, can be computed by a circuit of size $O(|C|k)$. Since, by construction of $k$, we have that $m \geq |C|^{\frac{2}{\epsilon}}$, a solution to {\sc $2^{\epsilon n}$-Hard} on input $1^m$ will be a truth table of length $m$ not computable by a circuit of size $m^{\epsilon} \geq |C|^2$, and thus a circuit of size $O(|C|k) = O(|C|\log |C|)$ would be a contradiction for all input lengths greater than some absolute constant. We construct such a circuit for any string in the range of $C^*$ as follows: let $y$ be a $2^kn$-bit string such that for some $x \in \{0,1\}^n$, $C^*(x)=y$. The circuit computing $y$ will have $x$ written as advice/constants, and will feed $x$ through $k$ copies of the circuit $C$ in series. We will split the $\lceil \log n \rceil + k$ input variables into a block of $k$ variables we call $i$, and a block of $\lceil \log n \rceil$ variables we call $j$. We then use $i$ to determine whether to apply $L$ or $R$ to the output of one of the copies of $C$ before feeding it into the next, to get some resulting string $x^i$, and then we use $j$ to index into the $j^{th}$ position of $x^i$, to get $y_{i,j}$. A diagram of this circuit is shown in Figure~\ref{fig:efficient_circuit}:

\begin{figure}[H]
    \centering
    \includegraphics[width=1\textwidth]{efficient_circuit_landscape.pdf}
    \caption{A succinct circuit whose truth table is $y$, for any $y$ in the range of $C^*$. Dotted boxes indicate the number of bits along a wire. Note that although $x$ is shown as an input in this diagram, for any given $y$ we fix a preimage $x$ as constants/advice, and so the only true inputs to this circuit are $i,j$.}
    \label{fig:efficient_circuit}
\end{figure}

To see that this circuit has size $O(|C|k)$, note that the subcircuits computing either $L$ or $R$ depending on a bit of $i$ can be computed easily with $O(n)$ gates over the basis $\{\land, \lor,\lnot\}$ (this is essentially a multiplexer), and also that the final subcircuit indexing into an $n$-bit string can be computed with $O(n)$ $\{\land, \lor,\lnot\}$ gates as well. Since $|C| \geq n$, and this circuit contains only $k$ copies of $C$ and the aforementioned subcircuits, plus the constants describing the string $x$ of length $n$, this circuit has $O(|C|k)$ size as claimed. 

 Thus, we now know that any solution to {\sc $2^{\epsilon n}$-Hard} on input $1^m$ will not be in the range of $C^*$, and by assumption {\sc $2^{\epsilon n}$-Hard} is total for sufficiently large input lengths so such a solution exists. It remains only to show that we can use a string outside the range of $C^*$, together with a $\cl{NP}$ oracle, to find a string outside the range of $C$. We proceed exactly as in the proof of \todo{reference stretching reduction}. Let $y$ be any string outside the range of $C^*$. Refer to Figure~\ref{fig:circuit_tree} which gives a diagram of a circuit computing $C^*$; at a layer $i \in [k]$ of this circuit, we have $2^i$ blocks of $n$ bits feeding into $2^i$ copies of $C$, and these copies of $C$ then output $2^{i+1}$ blocks of $n$ bits at the next layer. So working back from the output layer $k$, we can test if any consecutive $2n$-bit block of $y$ is outside of the range of $C$. If none of them are, then we find a preimage for all blocks, interpret this as the output of the previous layer, and continue our search from there. We follow this process all the way back to the input layer or until we find an empty pigeonhole of $C$. If we never find an empty pigeonhole of $C$, then this process will terminate at the input layer with a string $x$ such that $C^*(x) = y$, which is impossible by assumption, so at some point we must indeed find a string outside the range of $C$. Checking whether a particular string is an empty pigeonhole, or finding a preimage if it's not, can be accomplished with polynomially many callse to an $\cl{NP}$ oracle. We perform this test at most $2^k = poly(|C|)$ times (once for every copy of $C$ in the diagram in Figure~\ref{fig:circuit_tree}), so overall this process can be accomplished in polynomial time with an ${\cl NP}$ oracle.
\end{proof}

\todo{transition}

Finally, we turn our attention to a problem originally defined in \cite{KKMP}, which captures the principle: every graph with bounded degree and many vertices must have large diameter.
\begin{definition}[Remote Vertex]
Remote Vertex is the following problem. The input is a circuit $N: \B^n \to (\B^n)^m$, a parameter $\ell \in \mathbb{N}$, and some $v \in \B^n$, such that $n > \ell \lceil \log m \rceil$. The interpretation is that $N$ defines the outgoing neighborhood function of a directed graph on vertex set $\B^n$, where $N(u)$ lists the set of $\leq m$ neighbors of $u$ (possibly with repetition). A solution is some $u \in \B^n$ such that there is no directed path from $v$ to $u$ of length at most $\ell$.

In the Undirected Remote Vertex problem, the input is the same, but we allow as a solution any witness that $N$ fails to represent a undirected graph. In particular, we allow as a solution any pair $u,w \in \B^n$ such that $u \in N(w)$ but $w \notin N(u)$. 
\end{definition}
\todo{this problem captures som specail explicit construction problems, like lower bounds for words in groups; however we show it's just the same thing as avoid}

\begin{theorem}
Undirected Remote Vertex is equivalent to Avoid under many-one $\mathsf{FP}^{NP}$ reductions.
\end{theorem}
\begin{proof}
The reduction of Undirected Remote Vertex to Avoid is easy does not require an $\cl{NP}$ oracle. Say that the Undirected Remote Vertex instance has $2^n$ vertices, degree $m$, path length $\ell$, and start node $v \in \B^n$. We simply take our avoid instance $C: \B^{\ell \lceil \log m \rceil} \to \B^n$ to be the map which takes as input a description of a walk of length $\ell$ starting at $v$, and outputs the node that it reaches. Note that this reduction works just as well for the harder Remote Vertex problem.

We now prove the other direction. By \todo{prev}, we may assume the Avoid instance is of the form $C: \B^n \to \B^m$, $m \geq 2 n \log n$ (this is the only place we will use an $\cl{NP}$ oracle in our reduction). Consider a graph $G$ on vertex set $V = V_0 \times V_1$, with $V_0 = \mathbb{F}_2^n,  V_1 = \mathbb{F}_2^m$. A vertex $(x,y)$ has neighbors $(x + e_i, y)$ for each standard basis vector $e_i \in \mathbb{F}_2^n$, as well as the neighbor $(x, y+ C(x))$. This defines an undirected graph with degree at most $n + 1$ and $2^{n + m}$ vertices. For $\ell = 2n+2$, the instance $(G, 0^{n+m},\ell)$ defines a valid instance of of Undirected Remote Vertex provided $m + n  > \ell \lceil \log (n+1) \rceil $ which holds for our choice of parameters provided $n$ is sufficiently large. Notice that our instance defines an undirected graph by construction, so the only relevant solutions are vertices at distance $> \ell$ from $0^{n+m}$.
We claim that if $(x,y)$ is a valid undirected remote vertex solution, than $y \notin \rng(C)$; equivalently, every $(x,y)$ with $y \in \rng(C)$ is reachable from the origin in at most $2n + 1$ steps. Say that $y = C(z)$ In $n$ steps we may traverse from the start vertex $(0^n, 0^m)$ to $(z, 0^m)$ using bit flips, then in one step we may traverse from $(z, 0^m)$ to $(z, C(z) + 0^m) = (z, y)$, and finally in $n$ more steps we may traverse from $(z,y)$ to $(x, y)$ using bit flips again. 
\end{proof}
\section{Direct Reductions to Hard Truth Tables}
Ideally we could extend the main completeness result of the previous section \todo{reference} to work with polynomial time reductions, as opposed to ${\cl FP}^{\cl NP}$ reductions. However, the {\cl NP} oracle seems highly necessary for the proof techniques used above. Indeed, subsequent to the initial publication of this work, \cite{ILW} gave strong cryptographic evidence that the completeness result in Theorem~\ref{thm:apepp-completeness} requires an ${\cl NP}$ oracle, a result which we will discuss further \todo{somewhere?}. Despite this obstacle, we show that there are some important explicit construction problems, which are not \emph{obviously} about circuit complexity, which nonetheless can be reduced to the problem of finding truth tables of hard functions via {\cl FP} reductions. 

A simple way to phrase the following results is that any truth table with sufficiently large circuit complexity will necessarily satisfy a variety of other pseudorandom properties for which no explicit constructions are known, including: rigidity over $\mathbb{F}_2$, high space-bounded communication complexity, and high bit-probe complexity. To show this, we demonstrate that the failure of a string $x$ to possess any of these properties implies a smaller than worst case circuit for $x$. 

% \begin{definition}
% {\sc $\delta$-Quite Hard} is the following problem: given $1^N$, output an $N$-bit truth table with hardness $\frac{\delta N}{\log N}$
% \end{definition}
% This problem is total for sufficiently small $\delta$. We recall also the definition of {\sc $\epsilon$-Hard}, where we must construct a truth table of hardness $N^{\epsilon}$.
\todo{make sure the rigidity definition below matches}

\subsection{Rigidity}
We begin with the case of rigidity. We recall the definition of the search problem {\sc $(\epsilon,2)$-Rigid}, wherein given $1^N$ we are asked to produce an $N \times N$ matrix over $\mathbb{F}_2$ which requires changing $\geq \epsilon N^2$ entries to reduce its rank to $\epsilon N$. Any matrix achieving $(\epsilon,2)$ rigidity for an absolute constant $\epsilon >0$ (or even $\epsilon >> \frac{1}{\log \log N}$) would be sufficient to carry out Valiant's program \cite{Valiant}.

% \begin{definition}
% {\sc $\epsilon$-Rather Rigid} is the following search  problem: given $1^N$, construct an $N \times N$ matrix over $\mathbb{F}_2$ which is $(\epsilon N, \epsilon N^2)$-rigid.
% \end{definition}

% These parameters are the best possible up to constant factors, and in particular would be sufficient to carry out Valiant's program over $\mathbb{F}_2$.

\begin{theorem}~\label{thm:rigidity-circuit}
For any sufficiently small $\delta > 0$, there exists some $\epsilon > 0$ such that {\sc $\epsilon$-Rather Rigid} reduces in polynomial time to {\sc $\delta \frac{2^n}{n}$-Hard}.
\end{theorem}

\begin{proof}
To prove this, it suffices to show that for any $N \times N$ matrix $M$ which is not $(\epsilon N, \epsilon N^2)$-rigid, we can construct a Boolean circuit with $f(\epsilon)O(\frac{N^2}{\log N})$ gates which decides the value of $M_{i,j}$ given the $2\lceil \log N \rceil$-bit input $(i,j)$, for some function $f$ which approaches zero as $\epsilon$ approaches zero. This then implies that for any fixed $\delta$, an $N^2$-bit truth table requiring circuits of size $\frac{\delta N^2}{\log N}$ must be $(\epsilon N, \epsilon N^2)$-rigid  for some $\epsilon > 0$ which is a function only of $\delta$ (and otherwise determined by $f$ and the constants hidden in the $O(\cdot)$ term).

Say $M$ is not $(\epsilon N, \epsilon N^2)$-rigid. So there exists an $N \times \epsilon N$ matrix $L$, an $\epsilon N \times N$ matrix $R$, and an $N \times N$ matrix $S$ with at most $\epsilon N^2$ nonzero entries, such that $LR \oplus S = M$. We will construct a circuit allowing us to efficiently index $M$ which uses these matrices $L, R, S$ as advice. To encode $L$ and $R$, we can utilize the well-known theorem of Shannon that any truth table of length $N$ can be computed by a circuit of size $O(\frac{N}{\log N})$ \cite{Shannon}. Thus, $L$ can be specified as a list of $\epsilon N$ circuits, each of size $O(\frac{N}{\log N})$, where the $j^{th}$ circuit $L_j$ represents the $j^{th}$ column, and $L_j(i)$ computes $L_{i,j}$. The same can then be done for $R$ (indexing columns instead). To encode $S$, we employ a refinement of Shannon's result due to Lupanov \cite{Coin-Problem}, which tells us that for sufficiently large $N$, any truth table of length $N$ with at most $\epsilon N$ nonzero entries can be computed by circuit of size
\[
\frac{\log\binom{N}{\epsilon N}}{\log \log \binom{N}{\epsilon N}} + o\left(\frac{N}{\log N}\right) \leq H(\epsilon) O(\frac{N}{\log N})
\]
Where $H$ denotes the binary entropy function. Thus $S$ can be computed by a circuit of size $H(\epsilon) O(\frac{N^2}{\log N})$.

It remains to show that the additional circuitry we need to compute $M_{i,j}$ given $i, j$ and the encodings of $L,R,S$ does not increase things too much. By definition, we have that: 
\[M_{i,j} = \langle row_i(L), col_j(R) \rangle \oplus S_{i,j}\]
where the dot product is taken over $\mathbb{F}_2$. Now consider the circuit diagram shown in Figure~\ref{fig:rigid}:

\begin{figure}[H]
    \centering
    \includegraphics[width=0.3\textwidth]{Non_Rigid_Circuit.pdf}
    \caption{A small circuit for a non-rigid truth table}
    \label{fig:rigid}
\end{figure}
The subcircuit computing $\text{row}_i(L)$ is defined as follows: as described above, $L$ represents a $N \times \epsilon N$ matrix, specified as a list of $\epsilon N$ circuits of size $O(\frac{N}{\log N})$, each computing a column of $L$. $\text{row}_i(L)$ feeds $i$ into each of these circuits in parallel to get an $\epsilon N$-bit string giving value of the $i^{th}$ row of $L$. $\text{col}_j(R)$ is defined analogously. The circuit $S_{i,j}$ simply computes the $(i,j)^{th}$ entry of $S$ given $i,j$. Finally, the subcircuit $\oplus(X)$ computes the parity of its input string, the subcircuit $\land(X,Y)$ computes the bit-wise AND of two equal length strings, and the terminal gate computes the two-bit parity function. Given the previous equation relating the $(i,j)^{th}$ index of $M$ to the $(i,j)^{th}$ rows/columns/indices of $L, R, S$, it is straightforward to see that this circuit performs the necessary calculation.
It is clear that the $\oplus(X)$ and $\land(X,Y)$ can be implemented with a number of gates linear in their input size, which in this case is $\epsilon N$. From the analysis above, each of $\text{row}_i(L)$ and $\text{col}_j(R)$ can be implemented using $\epsilon O(\frac{N^2}{\log N})$ gates, and $S_{i,j}$ can be implemented using $H(\epsilon)  O(\frac{N^2}{\log N})$ gates. So overall, this circuit has size $(H(\epsilon) + \epsilon)O(\frac{N^2}{\log N})$. Since $H(\epsilon) + \epsilon$ approaches zero for decreasing  $\epsilon$, this implies that for any fixed $\delta > 0$, an $N^2$-bit truth table requiring circuits of size $\frac{\delta N^2}{\log N}$ must be $(\epsilon N, \epsilon N^2)$ rigid for some $\epsilon>0$ which is a function only of $\delta$ (and which is otherwise determined by the constants hidden in the $O(\cdot)$ expressions above).
\end{proof}

We thus conclude that if {\cl E} contains a language of circuit complexity $\Omega(\frac{2^n}{n})$, then there is a polynomial time construction of $(\Omega(n), \Omega(n^2))$-rigid matrices over $\mathbb{F}_2$.

\subsection{Space-Bounded Communication Complexity}

The class ${\cl PSPACE^{CC}}$ was defined originally in \cite{communication-classes} as a generalization of the class ${\cl PH^{CC}}$ to an unbounded alternation of quantifiers. We will not give this original definition, but rather a simplification due to \cite{space-communication}.

\begin{definition}
Let $f: \{0,1\}^n \times \{0,1\}^n \rightarrow \{0,1\}$. We say $f$ has a space-$s$ protocol if there is a deterministic protocol deciding $f$ of the following form. Alice receives $x \in \{0,1\}^n$, and Bob receives $y \in \{0,1\}^n$. There is an $s$-bit shared memory, and Alice and Bob alternate turns writing to this memory. On a player's turn, they can modify the contents of the $s$-bit shared memory as a function only of its previous contents and their private input $x$ or $y$, or decide to halt and output some $z \in \{0,1\}$ (this decision is also a function only of their input and the shared memory's previous state). This protocol is valid if $f(xy) = z$ for all $x,y$.
\end{definition}


By a result of Song \cite{space-communication}, we have that if $f  \in {\cl PSPACE^{CC}}$ then $f$ has a $poly(\log n)$ space protocol (Song uses a slightly different model where Alice and Bob have private $s$-bit tapes, but our model is at least as strong up to a doubling in space since sharing the tape only increases their ability to communicate). Due to a basic counting argument, most functions $f$ require $\Omega(n)$ space, and any such function must lie outside of $ {\cl PSPACE^{CC}}$. However, it has been a long standing open problem to give an explicit construction a communication matrix outside of even ${\cl PH^{CC}}$ \cite{space-communication}. We thus define the following search problem:

\begin{definition}
{\sc $\delta$-SPACE} is the following search problem. Given $1^N$, where $N=2^n$, output a communication matrix $\{0,1\}^n \times \{0,1\}^n \rightarrow \{0,1\}$ such that $f$ requires space-$\delta n$ communication protocols.
\end{definition}

\begin{theorem}\label{thm:space-circuit}
For any $\delta < \epsilon < \frac{1}{2}$, {\sc $\delta$-Space} reduces in polynomial time to {\sc $2^{(\frac{1}{2} + \epsilon)n}$-Hard}.
\end{theorem}

\begin{proof}
The proof follows the exact same strategy as above, constructing a smaller-than-worst-case circuit for any matrix with a $\delta n$-space protocol. We will omit the detailed construction of our circuit since it is very similar to that of the above proof, but give here a clear sketch of its essential structure. Let $f: \{0,1\}^n \times \{0,1\}^n \rightarrow \{0,1\}$, and let $N = 2^n$. We can represent a space $s$ protocol for $f$ as $2s + 4$ different $N \times 2^s$ binary matrices, with each of the first $2s$ matrices telling us how one of the two players will modify a certain cell of the shared memory as a function of its previous state and their input, and the last four matrices determining whether a certain player will halt or continue and the value they will output if they halt, again as a function of their input and the shared memory contents. For $s = \epsilon \log N$ (with $\epsilon < 1$), we can then take these matrices as advice to our circuit, for a total of $O(\epsilon N^{1 + \epsilon} \log N)$ bits of advice. We can then add circuitry to simulate $2^{s+1} = 2N^{\epsilon}$ steps of this protocol on a given input $x,y \in \{0,1\}^n \times \{0,1\}^n$, indexing into the advice matrices in order to determine what to do next using the same indexing constructions we had for the proof of Theorem~\ref{thm:rigidity-circuit}. Any matrix solvable by a space $s$ protocol will be solved by a protocol that halts after $2^{s+1}$ steps (otherwise it will loop forever), so simulating for $2^{s+1}$ steps suffices. Since each of the indexing operations can be implemented using a number of gates linear in $\epsilon N^{1 + \epsilon} \log N$ using the constructions from the proof of Theorem~\ref{thm:rigidity-circuit}, our overall circuit for $f$ will have size $O(2N^{\epsilon}\epsilon N^{1 + \epsilon} \log N) = O(N^{1 + 2\epsilon} \log N)$. Thus, for any $\epsilon < \frac{1}{2}$, a truth table of length $N^2$ requiring circuits of size $N^{1 + 2\epsilon}$ will require $\delta \log N$-space protocols for any $\delta < \epsilon$.
\end{proof}

\subsection{Bit Probe Lower Bounds}
In \cite{Elias-Flower}, Elias and Flower defined a broad model for studying the space/query complexity of data structure problems, known as the ``bit-probe model.'' 
\begin{definition}
Given two sets $D,Q$, a function $f: D \times Q \rightarrow \{0,1\}$ and an integer $b$, the bit-probe complexity of $f$ for space $b$, denoted $BC_b(f)$, is the minimum over all encodings $G: D \rightarrow \{0,1\}^b$ of the number of bits of $G(x)$ that need to be probed in order to determine $f(x,y)$ for the worst case $x \in D,y \in Q$, given access to $y$.
\end{definition}
In this general definition of a data structure problem, we think of $D$ as the set of all possible pieces of ``data'' we might wish to encode in our data structure, $b$ as the number of bits we can use to encode a piece of data, $Q$ as the set of queries we wish to answer about an encoded piece of data, and $f$ as telling us the correct answers to all data/query pairs. $BC_b(f)$ then tells us the minimum number of probes required by any space-$b$ data structure in order to answer every query correctly for every possible piece of data.

This model was investigated further by Miltersen \cite{Bro}, who showed, using a simple counting argument, that most problems require an infeasible amount of space/probes in this model, but pointed out that no explicit data structure problem is known to be infeasible in this sense. We thus define the following explicit construction problem:
\begin{definition}
{\sc $\delta$-Probe} is the following search problem: given $1^N$ where $N = 2^n$ output a truth table $f: \{0,1\}^n \times \{0,1\}^n \rightarrow \{0,1\}$ such that $BC_{2^{\delta n}}(f) > \delta n$.
\end{definition}

\begin{theorem}\label{thm:probe-circuit}
For any $\delta < 2\epsilon < 1$, {\sc $\delta$-Probe} reduces in polynomial time to {\sc $2^{(\frac{1}{2} + \epsilon)n}$-Hard}.
\end{theorem}
\begin{proof}

Say $f$ is a function $f: D \times Q \rightarrow \{0,1\}$ such that $BC_b(f) \leq k$. So in particular there is an encoding $G: D \rightarrow \{0,1\}^b$ such for any $x \in D$, $y \in Q$, given access to $y$ we can determine $f(x,y)$ using at most $k$ probes to $G(x)$. Let $H: Q \rightarrow [b]^k$ be the function which, given $y$, tells us which positions of $G(x)$ to query. Finally, let $\phi: \{0,1\}^k \times Q \rightarrow \{0,1\}$ be the function which determines $f(x,y)$ given the results of the probes and the value of $y$. We will now give a compact representation for $f$, and then use it to construct a circuit computing $f$.

Let $R$ be a $|D| \times b$ binary matrix whose rows are indexed by elements of $D$, where $R_{i,j}$ gives the value of the $j^{th}$ bit of $G(i)$. Let $S$ be a $|Q| \times k \lceil \log b \rceil$ matrix whose rows are indexed by elements of $Q$, such that the $i^{th}$ row of $S$ is $H(i) \in [b]^k$. Finally, let $Z$ be a $2^k \times |Q|$ matrix such that $Z_{i,j} = \phi(i,j)$.

Now, given $x \in D$, $y \in Q$, we can use $S,R,Z$ to determine $f(x,y)$ as follows. First find the $y^{th}$ row of $S$ to get $H(y)$, which is a list of $k$ indices in $[b]$. Next, find the $x^{th}$ row of $R$, which is the $b$-bit string $G(x)$. Then, probe the indices specified by $H(y)$ to get some $k$-bit string $w$, and finally output $Z_{w, y}$, which is precisely $\phi(w,y) = f(x,y)$.

As in the previous proofs, it can easily be shown that these indexing operations can be accomplished with circuits of size linear in $S,R,Z$. Therefore, $f$ has a circuit of size $O(|S| + |R| + |Z|) = O(|Q|k\log b + |D|b + 2^k|Q|)$. So for any function $D \times Q \rightarrow \{0,1\}^n$ requiring circuits of size $\omega(|Q|k\log b + |D|b + 2^k|Q|)$, we must have $BC_b(f) > k$. 

In particular, if we take $D = Q = \{0,1\}^n$, and for any fixed $\epsilon < 1$ we take $b= 2^{\epsilon n}$, $k = \epsilon n$, we get that for any function $f: \{0,1\}^{2n} \rightarrow \{0,1\}$ requiring circuits of size $\omega(2^{(1 +\epsilon) n})$, it must be the case that $BC_{2^{\epsilon n}}(f) > \epsilon n$, since 
\[
|Q|k\log b + |D|b + 2^k|Q| = 2^n \epsilon^2 n^2 + 2^n2^{\epsilon n} + 2^{\epsilon n}2^n = O(2^{(1+\epsilon) n})
\]
So if we take $\delta < \epsilon$, any truth table of length $N^2$ with hardness $N^{1+\epsilon}$, or in other words any solution to {\sc $2^{(\frac{1}{2} + \frac{\epsilon}{2})n}$-Hard} on input $1^{N^2}$, must satisfy $BC_{2^{\delta n}}(f) > \delta n$. 
\end{proof}

\section{A Brief Survey of Subsequent Developments}
Since the publication of \cite{Korten}, a significant body of work on Range Avoidance has emerged. We discuss here some of the key developments in the area by other researchers which will not otherwise be covered in the later chapters of this thesis. 
\begin{enumerate}
    \item Hardness of Avoid
    \item Algorithmic method
    \item CHR/Li
    \item Restricted variants of Avoid
\end{enumerate}
